% !TEX root = ../tcc.tex

\subsection{Trabalhos Correlatos}
==posso colocar os nossos aqui?==

\subsection{Definição do problema}
\label{sub:model}

Considere o sinal $x:\mathbb{Z}\to\mathbb{R}$ e sua Transformada de Fourier Discreta, do inglês \ac{DFT} $X(\Omega):\mathbb{R}\to\mathbb{C}$ dada por:

\begin{align}
    \label{eq:dft}
    X(\Omega) = \sum_{k = -\infty}^{\infty} x[k] e^{j\Omega k} , \Omega \in [-\pi,\pi]
\end{align}

O objetivo é determinar a frequência de máximo módulo no espectro de $x$, isto é, a frequência mais energética do sinal. Assim define-se $\Omega_m$ como:

\begin{align}
    \label{eq:omegam}
    \Omega_m \triangleq \underset{\Omega \in [-\pi,\pi]}{\arg \max} \lvert X(\Omega) \rvert
\end{align}

\subsection{Objetivos}
\subsubsection{Objetivo Geral}

Elaborar um algorítimo que consiga detectar a frequência de pico, conforme \ref{sub:model}, com complexidade idêntica (ou inferior) ao algorítimo estado da arte descrito em \ref{sub:stateoftheart}.

\subsubsection{Objetivos Específicos}

Os seguintes são os objetivos específicos:

\begin{enumerate}
    \item Mostrar as características matemáticas do método: provar convergência, verificar complexidade assintótica e tempo de convergência.
    \item Escrever a infraestrutura de teste para a validação do algorítimo (gravador de amostras e analisador de espectro).
    \item Contrastar a performance do algorítimo proposto com \ref{sub:stateoftheart}.
\end{enumerate}

\subsection{Estrutura do Texto}